\section*{Summary}

A charge amplifier integrates the charge induced on its electrode (\cref{sec:20.1}). By Ramo's theorem each pair induces $e$ times the fraction of the gap its carriers cross, and the source does the work $eV_{\mathrm{bias}}$ times the same fraction, so trapping lowers the read charge and the Luke heat together. With exponential survival over a drift length $\lambda_{\mathrm{d}}$, the Hecht relation gives the collected fraction of a face event, and for $\lambda_{\mathrm{d}} \gg d_{\mathrm{det}}$ each carrier of a bulk event loses a third of what the carrier of a face event loses. The Fano factor of the measured charge has three terms, the cascade Fano factor scaled by the collection, the spread of partial charges and the spread of the mean collection over event positions; the last grows with the deposit. At \SI{95}{\percent} mean collection in silicon the measured Fano factor of a \SI{100}{eV} recoil lies between 0.125 and 0.158 depending on how the charge is lost. Trapped charge distorts the field over a finite part of the crystal by a percent at \SI{1}{V/cm} once some $10^5$ charges per \si{cm^2} are trapped, which ordinary event rates reach between flashes. Carriers faster than sound emit Luke phonons by the acoustic Cherenkov process. The mean phonon energy is \SI{0.84}{meV} for silicon electrons at \SI{1}{V/cm} in a transport Monte Carlo. The published total rates are low by a factor $k_{\mathrm{s}}/k$, and the emission from traps at millikelvin is not known.

The prompt phonons down-convert by anharmonic decay, at a rate $\propto E_{\mathrm{ath}}^5$ for longitudinal phonons, and scatter on isotopes at a rate $\propto E_{\mathrm{ath}}^4$ (\cref{sec:20.2}). Their spectrum approaches a scaling form in $E_{\mathrm{ath}}/E_*(t)$, and a hard birth is forgotten once $E_*$ has fallen to about half its energy. In bulk silicon this takes \SI{6.8}{ns} for births near the band top and about \SI{1}{\micro s} for the transverse daughters of optical decay. Near an instrumented face the birth is not forgotten, and at \SI{1}{\micro s} that region is \SI{13.5}{\percent} of a \SI{4}{mm} crystal. Polarization is forgotten within a nanosecond. These times are continuum estimates, and they do not hold for sapphire, where isotope mixing is slower than decay.

At the interface with a sensor film the averaged transmission is 0.938 from silicon into aluminum and 0.368 back, with the elastic constants of aluminum at zero temperature, tied by the reciprocity $\mathfrak{n}_1\bar t_{12} = \mathfrak{n}_2\bar t_{21}$ (\cref{sec:20.3}); \SI{56}{\percent} of the phonon flux in the film lies beyond the transverse critical angle and escapes only after mixing. A phonon leaves a film of thickness $d_{\mathrm{f}}$ in $t_l = 4d_{\mathrm{f}}/(v_s\bar t_{\mathrm{esc}})$, \SI{0.19}{ns} (acoustic mismatch) or \SI{0.24}{ns} (diffuse mismatch) for \SI{60}{nm} of aluminum, comparable to the pair-breaking time. The fitted silicon-to-aluminum coefficient of 0.3 to 0.55 is an absorption probability, and it constrains the ratio of pair-breaking length to film thickness to about 5 to 23 and leaves the escape probability unconstrained.

In the film a phonon between $2\Delta_{\mathrm{BCS}}$ and $4\Delta_{\mathrm{BCS}}$ breaks one pair and gives the quasiparticle yield $Y_{\mathrm{qp}} = 2\Delta_{\mathrm{BCS}}/E_{\mathrm{ath}}$ (\cref{sec:20.4}). Above $4\Delta_{\mathrm{BCS}}$ the Kaplan cascade gives 0.590 without escape, with a maximum of 0.600 at $5.46\,\Delta_{\mathrm{BCS}}$, and 0.349 to 0.554 at $10\Delta_{\mathrm{BCS}}$ for $t_l/t_0$ from 0.5 to 8; the yield then falls across the band in which down-converted phonons arrive. Carriers that reach an aluminum electrode return 85 to \SI{92}{\percent} of their energy at or below $2\Delta_{\mathrm{BCS}}$, phonons that cross the crystal ballistically.

The transport is linear in the birth spectrum, so a detector is described by the collection kernels of its sensors, and every readout quantity is a linear functional of them: the collection contrasts $\varpi_O$ and $\varpi_{\mathrm{gap}}$, the position spread and the collection ratio $\mathfrak{r}_{\mathrm{coll}}$ (\cref{sec:20.5}). In the mean-field limit, for a \SI{4}{mm} silicon slab with a niobium-like and an aluminum film, Luke phonons below \SI{3}{meV} reach only aluminum and $\varpi_O = 0.39$, falling to 0.045 at a mean Luke energy of \SI{2.9}{meV}. With one film type, $\mathfrak{r}_{\mathrm{coll}} - 1 = s_{\mathrm{ng}} - (1 - \kappa_{\mathrm{sub}})s_{\mathrm{sub}}$: the near-gap gain is small in thick films, and the share of Luke energy below $2\Delta_{\mathrm{BCS}}$ lowers the ratio if it is lost. Phonons that a film lets out return to the crystal and continue the same cascade, so film escape changes the ratio by less than 0.002. With energy conserved the estimate lies between 0.81 and 1.02 over mean Luke energies of 0.75 to \SI{4.7}{meV}. Only a selective sink that deletes the phonons a film emits after its first pair break, as the lumped film assumes, raises one aluminum film to between 1.05 and 1.53, an upper envelope. The value 1.128 that would explain the silicon phonon-yield ratio of \cref{sec:19.5} without stored energy therefore lies above the transport estimate; only computed kernels for real detectors can exclude it, and the two-bias measurement of \cref{sec:19.5} decides.

A phonon readout tells the recoil types apart only through their pair numbers. Mode, frequency or angle selection of the prompt phonons cannot separate nuclear from electron recoils in bulk silicon and germanium: only the soft share of an electron recoil, about \SI{1.5}{\percent} of its energy at \SI{100}{eV}, survives the memory loss, and a position spread of the channel fraction above 0.008 hides it. The zero-field gap-energy channel, up to $0.28$ of the recoil energy at \SI{100}{eV} in a two-window reading, is unresolved: its size and sign depend on the release law and site of the gap energy, which are not known at millikelvin. A low-bias readout of the Luke share reads the ionization yield with phonons and needs a contrast $\varpi_O$ of 0.11 to 0.43 for a position spread of 0.02 to 0.10 at \SI{100}{eV}. As the Luke gain grows, the pair numbers of the two recoils converge at equal total phonon energy, and at \SI{100}{V} even an ideal phonon observable separates silicon recoils of \SI{220}{eV} by only 1.22 standard deviations at the central measured yield and Fano factor. This no-go at $E_R \le \SI{300}{eV}$ holds over the $1\sigma$ bands of the measured inputs and fails at $2\sigma$. These verdicts are inputs to the readout column of the guide to next targets (\cref{sec:21.7}); all of them assume transport that has not yet been validated against the relations of this chapter.

\Cref{tab:20-nuisances} collects the truncations and inputs of the chapter with their sizes. The measured charge and phonon signals, with these conditions, are what the exclusion limits of the next chapter start from.

\begin{table}[tbp]
  \caption{Nuisances of charge, phonon and quasiparticle readout. Type in the classes of \cref{sec:8.8}, normalization, tilt, threshold and new support, with ``delayed'' for an item that moves charge or energy in time over an exposure. Basis of the size: ``computed'' (exact or next-order result computed), ``model'' (computed in a model), ``est.''\ (parametric estimate), ``input spread'' (spread or limitation of a measured or computed input), ``unsized'' (no size known). The last column gives the section where the item is treated.}
  \label{tab:20-nuisances}
  \centering\scriptsize\setlength{\tabcolsep}{3pt}
  \begin{tabular}{@{}>{\raggedright\arraybackslash}p{0.26\linewidth}>{\raggedright\arraybackslash}p{0.12\linewidth}>{\raggedright\arraybackslash}p{0.35\linewidth}>{\raggedright\arraybackslash}p{0.09\linewidth}>{\raggedright\arraybackslash}p{0.10\linewidth}@{}}
    \toprule
    Truncation or input & Type & Size & Basis & Section \\
    \midrule
    Exponential survival of drifting carriers & normali\-zation & collected fraction 0.632 (face), 0.736 (uniform) at $\lambda_{\mathrm{d}} = d_{\mathrm{det}}$; saturating drift velocity caps it below one & computed; model & \cref{sec:20.1} \\
    Fano factor of the measured charge & normali\-zation (width) & 0.125--0.158 at \SI{100}{eV}, \SI{95}{\percent} collection; position term $1.2\times10^{-4}\,\bar N_{\mathrm{eh}}$ & computed & \cref{sec:20.1} \\
    Space charge between flashes & delayed & \SI{1}{\percent} field change at \SI{1}{V/cm} for 6.3--$12.6\times10^4$ trapped charges per \si{cm^2} & est. & \cref{sec:20.1} \\
    Athermal emission from traps & delayed & cross-sections not computed & unsized & \cref{sec:20.1} \\
    Luke emission rate and spectrum & normali\-zation (of $\mathfrak{r}_{\mathrm{coll}}$, $\varpi_O$) & transverse emission omitted; mean energy \SI{0.84}{meV} for electrons at \SI{1}{V/cm} & model; unsized & \cref{sec:20.1,sec:20.5} \\
    Memory loss of prompt phonons & tilt & \SI{6.8}{ns} (Si), \SI{47}{ns} (Ge) near the band top; about \SI{1}{\micro s} for the \SI{22}{meV} daughters in Si; fails for sapphire & est. & \cref{sec:20.2} \\
    Interface transmission model & normali\-zation & $t_l$ differs by 1.31 between acoustic and diffuse mismatch; anisotropy: $\bar t_{\mathrm{esc}} = 0.327$--0.405 & model & \cref{sec:20.3} \\
    Quasiparticle yield above $4\Delta_{\mathrm{BCS}}$ & normali\-zation & 0.590 without escape; 0.349--0.554 at $10\Delta_{\mathrm{BCS}}$ for $t_l/t_0 = 0.5$--8; kinetic equations higher by up to 0.04 & model & \cref{sec:20.4} \\
    Soft primary share of an electron recoil & tilt & $f_{\mathrm{s}} \simeq 0.017$; $(1 - Y_q)f_{\mathrm{s}} = 0.015$ at \SI{100}{eV} & est. & \cref{sec:20.2,sec:20.5} \\
    Collection ratio of Luke phonons & normali\-zation & 0.81--1.02 (energy conserved); selective sink up to 1.53; 1.128 not reached; detector kernels not computed & model & \cref{sec:20.5} \\
    Collection contrast of Luke phonons & normali\-zation & 0.39 below about \SI{1.5}{meV}, 0.045 at \SI{2.9}{meV}; the low-bias route needs 0.11--0.43 & model & \cref{sec:20.5} \\
    Position spread of the channel fraction & normali\-zation & above 0.008 it closes the primary route & unsized & \cref{sec:20.5} \\
    Release law and site of the gap energy & tilt; delayed & two-window difference up to $0.28(1 - f_{\mathrm{ret}})$ at \SI{100}{eV}; sign not known & unsized & \cref{sec:20.5} \\
    Inputs of the high-voltage ceiling & threshold & \SI{220}{eV}, \SI{100}{V}: 1.22 (central), 1.50 (mean charge); 2.46 at \SI{100}{eV} with $2\sigma$ bands & input spread & \cref{sec:20.5} \\
    \bottomrule
  \end{tabular}
\end{table}
